%% ma: L11-B15-0076
%% lop: 11
%% chuong: 5
%% bai: 15
%% dang: 11.15.3
%% loai: DS
%% muc: [TH, TH, VD, VD]
%% dap_an: "SĐSĐ"
%% nguon: "tai-lieu-hoc-tap-mon-toan-lop-11-theo-chuong-trinh-giao-duc-pho-thong-2018.pdf (Lớp Toán Thầy Duy, Phạm Lê Duy, THPT Chu Văn An), Chương 3 Bài 1, Phần 2, Câu 6 (tr.337)"
%% kien_thuc: "Lớp 11 Bài 15; Bài 7 (cấp số nhân)"
%% trang_thai: chinh-thuc
%% ghi_chu: "Đã duyệt 03/11/2026."
\begin{ex}
Biết $\lim\dfrac{2n^2-n+4}{an^2+n+3}=2$ và $\lim\dfrac{3^n+4^{n+1}}{4^n+3}=b$.
\choiceTF
{Giá trị của $a=2$.}
{\True Giá trị của $b=4$.}
{$2a-b=0$.}
{\True Ba số $a,b,16$ lập thành một cấp số nhân.}
\loigiai{\sloppy $\lim\dfrac{2n^2-n+4}{an^2+n+3}=\dfrac2a=2\Rightarrow a=1$.
\[\lim\dfrac{3^n+4^{n+1}}{4^n+3}=\lim\dfrac{\left(\frac34\right)^n+4}{1+3\left(\frac14\right)^n}=4\Rightarrow b=4.\]
a) $a=1\ne2$. Sai.
b) $b=4$. Đúng.
c) $2a-b=2-4=-2\ne0$. Sai.
d) $1;4;16$ là cấp số nhân công bội $4$. Đúng.}
\end{ex}
